
\begin{claim}{AC.31}{Proved counterexample to a scalar-Frobenius reduction}
Let
\[
 C/\mathbf F_{16}:\quad y^2+y=x^5,\qquad P=\infty.
\]
Arithmetic Frobenius on \(J_C\) is the scalar \([-4]\).
Nevertheless there is a constant \(P\)-split cyclic cover of degree
thirteen over \(\mathbf F_{4096}\), its first possible constant-deck
field, whose new factor is ordinary of two-rank twelve and has odd
relative class number. More precisely its twisted polynomials are
\[
 L_P(\chi^i,T)
 =1+\frac{139-13\left(\frac{i}{13}\right)\sqrt{13}}2\,T
       +4096T^2,\qquad 1\le i\le12,
\]
after choosing the sign of \(\sqrt{13}\) consistently with the
quadratic Gauss sum. Consequently
\[
 \frac{P_D(T)}{P_C(T)}
 =\bigl(1+139T+12473T^2+569344T^3+16777216T^4\bigr)^6
\]
and
\[
 \frac{h_{\mathbf F_{4096}}(D)}{h_{\mathbf F_{4096}}(C)}
                         =17359173^6.
\]
Thus scalar Jacobian Frobenius, even at a rational hyperelliptic
ramification point, does not reduce the rational criterion to
vanishing of the geometric new \(p\)-rank.
\end{claim}
\begin{proof}
The exact counts \(\#C(\mathbf F_{16})=33\) and
\(\#C(\mathbf F_{256})=193\) give Frobenius polynomial
\((X+4)^4\). Frobenius acts semisimply, so it equals \([-4]\)
as a Jacobian endomorphism; equivalently apply
\cite[Theorem~2(d), p.~140]{AC_Tate1966}.
The following computation is independent of the assertion that all
thirteenth-degree Pryms of this curve are ordinary: it proves
ordinariness for the displayed cover itself.

Use binary polynomial notation in
\[
 K=\mathbf F_2[t]/(t^{12}+t^3+1),\qquad \#K=4096.
\]
The reduced divisor class
\[
 a=(u,v),\qquad
 u=x^2+3377x+483,\qquad v=3144x+2626
\]
has exact order thirteen. Its sixteenth-power conjugate is
\([-4]a\), and its \(4096\)-power conjugate is \(a\).
These are checked by exact Mumford arithmetic. Since
\((-4)\bmod13\) has order three and \((-4)^2=16\), all nonzero
geometric thirteenth-degree quotient characters have first
constant-deck field \(\mathbf F_{16^3}\). Also
\(4096\equiv1\bmod13\), so the displayed \(a\) has the constant-deck
descent required in AC.27.

For an independent direct character count, write \(u^{13}=M\).
Lift \(v\) to the unique solution of \(w^2+w=x^5\bmod M\) with
\(w\equiv v\bmod u\). Exact linear algebra gives polynomials
\(A,B\), of degrees thirteen and ten, with \(A\) monic and
\[
 A^2+AB+B^2x^5=u^{13},\qquad A+Bw\equiv0\bmod u^{13}.
\]
Their coefficient lists, in increasing degree, are
\[
\begin{aligned}
 A={}&(1334,2010,3185,1642,2095,2911,77,\\
     &\hspace{12mm}567,182,1378,3843,297,2491,1),\\
 B={}&(1250,517,799,137,81,3560,2285,3410,2702,2433,1191).
\end{aligned}
\]
Hence \(F=A+By\) has divisor \(13D_a-26P\), where
\(D_a\) is the effective degree-two divisor representing \(a\).
The norm identity and the vanishing condition distinguish its
two zeroes from their hyperelliptic conjugates.
The cover is the normalization of \(z^{13}=F\).
It is unramified because every valuation of \(F\) is divisible
by thirteen. It is connected because \(a\ne0\), and it is
\(P\)-split because the leading local unit of \(F\) at \(P\)
is one. Its deck group is constant since \(\mu_{13}\subset K\).

Let the primitive field generator be the binary element \(3\);
use its associated primitive thirteenth root to label the
character exponents. At a point where \(F\ne0\), the exponent is
\(\log_3F\bmod13\). At a zero \(R=(r,y)\), the norm identity
shows that the local unit has character exponent
\(-\log_3B(r)\bmod13\): the other factor of \(u^{13}\) contributes
a thirteenth power. Infinity contributes exponent zero.
Enumerating all \(4353\) rational points, with these local units
at the two zeroes, gives
\[
 (N_0,\ldots,N_{12})
  =(399,323,336,323,323,336,336,336,336,323,323,336,323).
\]
Thus the nonzero quadratic residues modulo thirteen have count
\(323\), and the nonresidues have count \(336\).
The quadratic-period identity
\(\sum_{r\ {\rm square}}\zeta_{13}^r=(-1+\sqrt{13})/2\)
gives
\[
 \sum_R\chi([R-P])
 =399+323\frac{-1+\sqrt{13}}2
       +336\frac{-1-\sqrt{13}}2
 =\frac{139-13\sqrt{13}}2.
\]
AC.4, with \(K_C\sim2P\), now gives every displayed
character polynomial and their product, including the exact
relative class number.

As a second check, all twelve Cartier scalars of AC.27 belong to
\(\{72,73\}\subset K\), with six of each value.
Both satisfy \(X^2+X+1=0\). Thus none is zero or one, and
\[
 \overline{P_D(T)/P_C(T)}
   =(1+T+T^2)^6=1+T^2+T^6+T^{10}+T^{12}.
\]
The degree twelve proves geometric new two-rank twelve, while
the value one at \(T=1\) proves the rational two-unit condition.
The program \path{scalar_frobenius_counterexample.py} executes
both the Cartier test and the direct Kummer-point test and
requires their complete polynomials modulo two to agree.

The failure of the proposed general reduction is structural.
On its rank-one \(\mathbf F_{4096}\)-module, the nonconstant
deck Frobenius over \(\mathbf F_{16}\) is semilinear for
\(x\mapsto x^{2^8}\), since \(-4\equiv2^8\bmod13\).
Its cube is a norm to \(\mathbf F_{16}^{\times}\), which need
not be one. Hyperelliptic reflection further restricts the scalar
to \(\mathbf F_4^\times\); it still leaves precisely the two
nontrivial values that occur here. Passing from the witness field
to its cubic constant extension raises these values to one and
destroys the rational criterion, giving an explicit instance of
the unsafe-extension warning in AC.6.
\end{proof}

\begin{claim}{AC.32}{Proved; a simultaneous genus-two counting criterion}
Let \(C/k\) have genus two, with \(k\) a finite field of
characteristic \(p\), and let \(P\in C(k)\).
Let \(\ell\ne p\) be an odd prime such that \(p\) is a primitive
root modulo \(\ell\). Set
\[
 V=J_C(k)/\ell J_C(k),\qquad H=\ell J_C(k),\qquad N=\#C(k),
\]
and assume \(V\ne0\). If either
\[
                  |V|>N\qquad\hbox{or}\qquad N\equiv-1\bmod p,
\]
and
\[
                         (p-1)(N-1)>2(|H|-1),                 \tag{AC.16}
\]
then some surjection \(J_C(k)\to\mathbf Z/\ell\) gives a
\(P\)-split L1 witness over \(k\).

One may replace (AC.16) by the sharper occupancy-dependent
condition
\[
 N_0(N_0-1)+(p-1)(N-N_0)>2(|H|-1),\qquad
 N_0=\#\{R\in C(k):[R-P]\in H\}.
\]
The subgroup \(H\) is the image of multiplication by \(\ell\),
not the much larger kernel of one chosen cyclic character.
No universal assertion about the existence of such a prime and
field is implicit in this conditional criterion.
\end{claim}
\begin{proof}
Suppose every prime cyclic quotient fails the rational criterion.
For each quotient, the nontrivial character evaluations of
\(\Theta=[0]+\sum_{R\in C(k)}[[R-P]\bmod H]\) are permuted
transitively by \(p\)-power Frobenius, because \(p\) is primitive
modulo \(\ell\). Failure means at least one is zero; hence all
are zero. Every nontrivial character of \(V\) factors through
one such quotient, so Fourier inversion in \(\mathbf F_p[V]\)
gives
\[
                 \Theta=c\sum_{b\in V}[b],
                 \qquad c=(N+1)/|V|\in\mathbf F_p.
\]
The denominator is invertible since \(\ell\ne p\).
If \(|V|>N\), some nonzero coset is empty and forces \(c=0\).
The alternative congruence \(N\equiv-1\bmod p\) also forces
\(c=0\) directly. Thus
\[
              N_0\equiv-1\bmod p,\qquad
              N_b\equiv0\bmod p\quad(b\ne0).
\]

For every nonzero Jacobian class \(\delta\), there are at most
two ordered pairs of geometric points \((R,S)\) with
\([R-S]=\delta\).
Indeed \(R-S=R'-S'\) implies
\(R+S'\sim R'+S\). If these effective degree-two divisors are
different, both belong to the unique hyperelliptic pencil;
the only possible second representation is
\((R',S')=(\iota S,\iota R)\).
Consequently the number of ordered distinct pairs in the same
\(H\)-coset is bounded above by \(2(|H|-1)\).
The occupancy congruences give the lower bound
\[
 \sum_bN_b(N_b-1)
   \ge N_0(N_0-1)+(p-1)(N-N_0).
\]
Since \(P\) lies in the zero coset, \(N_0\ge1\); its congruence
improves this to \(N_0\ge p-1\). Therefore
\[
\begin{aligned}
 N_0(N_0-1)+(p-1)(N-N_0)
   &=(p-1)(N-1)+(N_0-1)(N_0-(p-1))\\
   &\ge(p-1)(N-1).
\end{aligned}
\]
Either stated strict inequality contradicts the upper bound.
Thus at least one cyclic quotient is a witness.
The primitive-root hypothesis was used at the first step and
cannot be dropped by this argument.
\end{proof}

